\subsection{Prototype Query Interface}
\label{subsec:flask-prototype}

To demonstrate the practical queryability of the knowledge graph beyond
command-line SPARQL, a lightweight web-based query interface was developed
as a proof-of-concept. The prototype, implemented as a Flask application
backed by RDFLib, provides two query modalities: a structured form with
dropdown and checkbox filters for Annex~III domain, fundamental rights,
and system pattern; and a raw SPARQL textarea for advanced users.

The structured query mode translates user selections into SPARQL queries
dynamically, abstracting the graph's namespace and predicate structure
from the end user. A deployer conducting a Fundamental Rights Impact
Assessment under Article~27 can, for example, select ``Employment'' and
``Non-discrimination'' to retrieve all precedent incidents where
employment-domain AI systems implicated non-discrimination rights. Results
are displayed in a tabular format showing the incident title, source type,
domain, rights, harms, and system pattern for each matching record.

This prototype serves as evidence toward the Article~27(5) goal of
automated tooling to support FRIA compliance. While the current
implementation is minimal --- it does not include authentication, pagination,
or integration with external incident databases --- it demonstrates that the
semantic representation developed in this work is directly amenable to
deployment as a queryable evidence base. The transition from a static
knowledge graph to an interactive retrieval tool requires no schema
modification, confirming that the RDF/Turtle representation is both
machine-readable and operationally useful. Future work could extend this
prototype to incorporate federated SPARQL queries across multiple
knowledge graphs, or to integrate with existing AI risk management
platforms.
